\section{An invariant answers what cannot change}\label{ch05:sec:invariant}
Many problems describe a process: there is a collection of possible \emph{states}, and there are \emph{legal moves}, each of which turns one state into another. An \emph{invariant} of such a process is a quantity or property that every legal move leaves unchanged.

Here is an example. A row contains $n$ signs, each a plus or a minus, where $n\ge2$. A legal move picks two different positions in the row and flips the sign in each: a plus becomes a minus, and a minus becomes a plus. Can we start from a row of all plus signs and reach a row with exactly one minus sign?

To find an invariant, replace each plus sign by the number $1$ and each minus sign by $-1$, and multiply all the numbers in the row together. Call the result the \emph{product of the row}. Flipping one sign multiplies that entry by $-1$, and therefore multiplies the whole product by $-1$. A legal move flips two signs, so it multiplies the product by $(-1)(-1)=1$: the product never changes. For instance, the row $+\,-\,+\,+$ has product $1\cdot(-1)\cdot1\cdot1=-1$. Flipping the first two signs gives $-\,+\,+\,+$, whose product is still $-1$.

A row of all plus signs has product $1$. Induction on the number of moves shows that every row reachable from it also has product $1$. After $0$ moves the product is $1$. If the product is $1$ after $k$ moves, the next move leaves it unchanged, so it is still $1$ after $k+1$ moves. A row with exactly one minus sign has product $-1$, so it can never be reached, however many moves we make.

This is an impossibility proof, and it does not examine the possible sequences of moves one by one. That would be hopeless: moves can be repeated, so there are infinitely many sequences. Instead the proof finds a barrier that every sequence must respect.

The product also has a simple meaning. It equals $-1$ raised to the number of minus signs, so it is $1$ exactly when the number of minus signs is even. Seen this way, a legal move changes the number of minus signs by $+2$ (two plus signs flipped), by $-2$ (two minus signs flipped) or by $0$ (one of each), so that number stays even. In this example the invariant tells the whole story. Every row with an even number of minus signs \emph{can} be reached from the all-plus row: choose the positions that should hold minus signs, pair them up, and flip one pair at a time. In general, though, an invariant can only show that a state is unreachable. To show that a state can be reached, we must exhibit moves that actually get there.

An invariant also says nothing about whether a process stops. The sign-flipping process, for instance, can run forever: flip the first two signs, flip them back, and repeat. Every row along the way has product $1$, yet nothing forces the process to end. To show that a process must stop, we need a separate argument that it makes \emph{progress}. The next section shows both kinds of argument working together.
